\documentclass[11pt]{article}
\usepackage[T1]{fontenc}
\usepackage{lmodern}
\usepackage{amsmath,amssymb,amsthm}
\usepackage[margin=1in]{geometry}
\usepackage{microtype}
\usepackage[colorlinks=true,linkcolor=blue,citecolor=blue,urlcolor=blue]{hyperref}
\hypersetup{pdftitle={Norm inflation in homogeneous H one-quarter for the defocusing quintic Gabitov--Turitsyn equation},pdfauthor={Kaiwen Jin}}
\newcommand{\R}{\mathbb R}
\numberwithin{equation}{section}
\newtheorem{theorem}{Theorem}[section]
\newtheorem{lemma}{Lemma}[section]
\newtheorem{corollary}{Corollary}[section]
\title{Norm inflation in $\dot H^{1/4}(\R^3)$ for the defocusing quintic Gabitov--Turitsyn equation}
\author{Kaiwen Jin\thanks{Contact: \href{mailto:yui.ui.siki@gmail.com}{\texttt{yui.ui.siki@gmail.com}}.}}
\date{}
\begin{document}
\maketitle
\begin{abstract}
We prove norm inflation in the homogeneous Sobolev space $\dot H^{1/4}(\R^3)$
for the defocusing quintic Gabitov--Turitsyn equation with averaging interval
$[0,1]$. A family of Schwartz initial data has arbitrarily small initial norm,
while its global smooth solutions attain arbitrarily large norms at arbitrarily
small negative times. The proof combines a finite-interval virial identity with
a lower bound for the averaged nonlinear forcing obtained from localized
low-frequency mass. A rescaling keeps the variance bounded over a long time
interval and forces growth of the fractional Sobolev norm.
\end{abstract}

\section{Introduction}\label{sec:intro}

The Gabitov--Turitsyn equation is a nonlinear Schr\"odinger equation whose
nonlinearity is averaged along the free Schr\"odinger flow. We consider the
three-dimensional defocusing quintic equation
\begin{equation}\label{eq:GT}
 iu_t-\Delta u+Q_1(u)=0,\qquad
 Q_1(f)=\int_0^1T_{-\sigma}(|T_\sigma f|^4T_\sigma f)\,d\sigma,
 \qquad T_\sigma=e^{i\sigma\Delta}.
\end{equation}
The averaging interval is fixed throughout. For this sign and power, global
well-posedness in the energy space with a finite second spatial moment follows
from Kawakami and Murphy~\cite[Proposition~5.1]{KM}. This global control leaves
open the behavior of weaker Sobolev norms for initial data whose energy is large.

Kowalski~\cite[Theorem~1.2]{K} proves failure of $C^5$ regularity of the solution
map below the mass-scaling index $s_m=1$ for this model. Norm inflation asks
for a stronger instability. The integrated-scaling index is $s_i=1/2$;
the range $0<s<s_i$ is left open in the discussion
of~\cite[Section~6.4]{K}. We resolve the fixed exponent $s=1/4$.

\begin{theorem}\label{thm:inflation}
For every $\varepsilon>0$, there are $u_0\in\mathcal S(\R^3)$ and a global
smooth solution $u$ of~\eqref{eq:GT}, with $u(0)=u_0$, such that for some
$t_\varepsilon\in(-\varepsilon,0)$,
\[
 \|u_0\|_{\dot H^{1/4}(\R^3)}<\varepsilon,
 \qquad \|u(t_\varepsilon)\|_{\dot H^{1/4}(\R^3)}>\varepsilon^{-1}.
\]
\end{theorem}

The time in the theorem is negative and nonzero. The initial data are Gaussians
with growing amplitude and shrinking width. Their $H^1$ norms need not be
small, so the conclusion is compatible with global energy-space control.

Our argument uses the defocusing virial approach of
Kowalski~\cite[Proposition~1.5 and Section~6.1]{K}, keeping the finite endpoint
of the rescaled averaging interval. Growth of the kinetic energy alone does
not give a lower bound for a fractional Sobolev norm: a small amount of mass
at very high frequency can carry the kinetic energy. Instead, a bounded
variance and a small fractional norm keep a fixed amount of mass in a fixed
ball for a short free evolution. This yields a lower bound for a positive
term in the variance identity. After rescaling, its time integral is bounded
independently of the long observation interval. The resulting incompatibility
forces growth in $\dot H^{1/4}$.

Section~\ref{sec:estimates} proves the localization and virial estimates and
justifies the regularity used in them. Section~\ref{sec:proof} gives the
rescaling and the quantitative Gaussian construction. We make no claim here
about the other positive exponents below $1/2$ or about inflation at positive
times.

We use the unitary Fourier transform
$\widehat f(\xi)=(2\pi)^{-3/2}\int_{\R^3}e^{-ix\cdot\xi}f(x)\,dx$ and
\[
 \|f\|_{\dot H^s}^2=\int_{\R^3}|\xi|^{2s}|\widehat f(\xi)|^2\,d\xi.
\]
The spaces $H^k$ are the inhomogeneous Sobolev spaces, $\|\cdot\|_p$ denotes
the $L^p(\R^3)$ norm, and $\langle f,h\rangle=\int\bar f h$.
All unspecified spatial integrals are over $\R^3$. Constants in the estimates
may depend on a displayed derivative order or on fixed data, but the constants
in the norm-inflation construction are independent of its parameter $g$.

\section{Two estimates for the averaged equation}\label{sec:estimates}

We first record the regularity needed for the virial calculation. Write
$\Sigma=\{f\in H^1(\R^3):xf\in L^2(\R^3)\}$.

\begin{lemma}[Smooth Gaussian solutions]\label{lem:smooth}
For $A,L>0$ and $u_0(x)=A\exp(-|x|^2/(4L^2))$, equation~\eqref{eq:GT}
has a global smooth solution $u$. On every compact time interval $I$,
\[
 u,\ xu,\ |x|^2u\in C(I;H^k(\R^3))\qquad(k=0,1,2,\ldots).
\]
\end{lemma}

\begin{proof}
Kawakami and Murphy~\cite[Proposition~5.1]{KM} prove global well-posedness
in $\Sigma$, with a uniform $H^1$ bound, for the defocusing equation
$iv_t+\Delta v=Q_1(v)$ when $d=3$ and the nonlinearity is $|v|^4v$. The change
$u(t)=v(-t)$ gives~\eqref{eq:GT} with the same averaging interval.

Here are persistence estimates for the smooth weighted data used below.
Within this proof, $D^k$ denotes the family of order-$k$ derivatives.
For each integer $k\ge1$,
\begin{equation}\label{eq:tame}
 \|Q_1(f)\|_{H^k}\le C_k\|f\|_{H^1}^4\|f\|_{H^k}.
\end{equation}
Indeed, the free Strichartz estimate~\cite[Theorem~1.2]{KT} for
$(q,r)=(5,30/11)$ and Sobolev embedding give
\[
 \|T_\sigma D^kf\|_{L^5_\sigma L^{30/11}_x}\lesssim_k\|f\|_{H^k},
 \qquad
 \|T_\sigma f\|_{L^5_\sigma L^{30}_x}\lesssim\|\nabla f\|_2,
 \qquad 0\le\sigma\le1.
\]
In a Leibniz term, let $k_1+\cdots+k_5=k$ be the derivative orders on its
five factors. Integer Gagliardo--Nirenberg interpolation~\cite[p.~125]{N}
gives, with $h=T_\sigma f$,
\[
 \|D^{k_j}h\|_{p_j}
 \lesssim_k\|h\|_{30}^{1-k_j/k}\|D^kh\|_{30/11}^{k_j/k},
 \qquad \frac1{p_j}=\frac{1+10k_j/k}{30}.
\]
The cases $k_j=0,k$ are identities. The spatial reciprocals sum to $1/2$;
five $L^5_\sigma$ norms give an $L^1_\sigma$ product. Unitarity of
$T_{-\sigma}$ proves the highest-derivative part of~\eqref{eq:tame}.
For its $L^2$ part use one $L^{30/11}_x$ factor and four $L^{30}_x$ factors.
Summing the lower derivatives proves~\eqref{eq:tame}. The admissibility
condition is $2/5+3/(30/11)=3/2$.

For the weights, set $J_\sigma=x+2i\sigma\nabla$ and $F(z)=z^3\bar z^2$.
The identities
\[
 J_\sigma T_\sigma f=T_\sigma(xf),\qquad
 J_\sigma F(h)=3h^2\bar h^2J_\sigma h
                 -2h^3\bar h\,\overline{J_\sigma h}
\]
and the $H^k$ algebra property, $k\ge2$, yield
\begin{equation}\label{eq:firstweight}
 \|xQ_1(f)\|_{H^k}\le C_k\|f\|_{H^k}^4\|xf\|_{H^k}.
\end{equation}
Expanding $\sum_{j=1}^3(x_j+2i\sigma\partial_j)^2F(T_\sigma f)$ also gives
\begin{equation}\label{eq:secondweight}
 \||x|^2Q_1(f)\|_{H^k}\le C_{k,R}(1+\||x|^2f\|_{H^k})
 \quad\text{if }\|f\|_{H^{k+2}}+\|xf\|_{H^{k+1}}\le R.
\end{equation}
To see the dependence in this estimate, the free position commutator bounds
$\||x|^2T_\sigma f\|_{H^k}$ by a constant times
$\||x|^2f\|_{H^k}+\|xf\|_{H^{k+1}}+\|f\|_{H^{k+2}}$ for $0\le\sigma\le1$.
The terms with $x\nabla F$ contain one weighted factor, and those with
$\Delta F$ contain at most two further derivatives.

Local contraction can be run in the sum-norm space
$H^{k+2}\cap\{xf\in H^{k+1},\ |x|^2f\in H^k\}$; the corresponding
polynomial difference estimates give local Lipschitz continuity.
This solution lies in $\Sigma$ and is identified with the global solution above.
On a compact interval,~\eqref{eq:tame} and the $H^1$ bound first control all
higher Sobolev norms by Gronwall. The commutators
\[
 [\Delta,x]=2\nabla,\qquad [\Delta,|x|^2]=4x\cdot\nabla+6
\]
then let~\eqref{eq:firstweight} control the first weight, followed by
\eqref{eq:secondweight} for the second. These bounds give continuation in the
same weighted space. Sobolev embedding and the equation give smoothness in
space and time.
\end{proof}

The first estimate controls the positive term that will appear in the variance
identity. A small fractional norm leaves most of the mass below a suitable
frequency cutoff; its position norm stays bounded during a short free evolution.

\begin{lemma}[Free mass localization]\label{lem:localization}
Fix $M,Y>0$ and $0<s<1$. There is $c=c(M,Y,s)>0$ such that, for $S\ge1$ and
$f\in H^1(\R^3)$ satisfying $\|f\|_2^2=M$ and $\|xf\|_2\le Y$,
\begin{equation}\label{eq:staticlower}
 W_S(f):=\int_0^S\sigma\|T_\sigma f\|_6^6\,d\sigma
 \ge c(1+\|f\|_{\dot H^s})^{-2/s}.
\end{equation}
\end{lemma}

\begin{proof}
Choose $\chi\in C_c^\infty(\R^3)$ with $0\le\chi\le1$, $\chi=1$ on
$\{|\xi|\le1\}$, and $\chi=0$ on $\{|\xi|\ge2\}$. Let $P_N$ be the Fourier
multiplier with symbol $\chi(\xi/N)$, and set $f_N=P_Nf$. Put
$B=\|f\|_{\dot H^s}$ and
\[
 N=\max\{1,[16(1+B)/\sqrt M]^{1/s}\}.
\]
Plancherel gives
\[
 \|f-f_N\|_2\le N^{-s}B\le\frac{\sqrt M}{16},
 \qquad \|f_N\|_2\ge\frac{15\sqrt M}{16}.
\]
Differentiating the cutoff in frequency gives
\[
 \|xf_N\|_2\le Y+C_\chi\frac{\sqrt M}{N},
 \qquad \|\nabla f_N\|_2\le2N\sqrt M.
\]
Since $xT_\sigma=T_\sigma(x-2i\sigma\nabla)$, for $0\le\sigma\le1/N$,
\[
 \|xT_\sigma f_N\|_2\le Y+(C_\chi+4)\sqrt M=:Y_*.
\]
Fix $R\ge1$ so that $Y_*/R\le\sqrt M/8$. The mass outside the ball then
satisfies $\|\mathbf1_{|x|>R}T_\sigma f_N\|_2\le\sqrt M/8$. Subtracting
this tail and the high-frequency remainder gives
\[
 \|\mathbf1_{|x|\le R}T_\sigma f\|_2
 \ge(15/16-1/8-1/16)\sqrt M\ge\sqrt M/2.
\]
H\"older on the ball $B_R=\{|x|\le R\}$ yields
\[
 \|T_\sigma f\|_6^6\ge\frac{(M/4)^3}{|B_R|^2}=:c_*.
\]
The interval $[0,1/N]$ lies in $[0,S]$, so $W_S(f)\ge c_*/(2N^2)$.
Finally $N\le C_{M,s}(1+B)^{1/s}$ proves~\eqref{eq:staticlower}.
The cutoff is fixed once; all constants are independent of $f$ and $S$.
\end{proof}

We next keep track of the finite upper endpoint in the averaged nonlinearity.
For $S>0$ write
\[
 Q_S(f)=\int_0^S T_{-\sigma}(|T_\sigma f|^4T_\sigma f)\,d\sigma.
\]
For a solution of
\begin{equation}\label{eq:scaled}
 iw_\tau-\beta\Delta w+Q_S(w)=0,\qquad \beta>0,
\end{equation}
let
\[
 V=\|xw\|_2^2,\qquad
 D=\operatorname{Im}\int_{\R^3}\bar w\,x\cdot\nabla w\,dx,
 \qquad G(\sigma,\tau)=\|T_\sigma w(\tau)\|_6^6,
\]
\[
 P(\tau)=\int_0^S G(\sigma,\tau)\,d\sigma,
 \qquad W(\tau)=W_S(w(\tau)).
\]
Here $V$ is the second spatial moment and $D$ is its dilation moment.
The mass and positive energy are
\begin{equation}\label{eq:energy}
 M=\|w\|_2^2,\qquad
 E=\frac\beta2\|\nabla w\|_2^2+\frac16P.
\end{equation}

\begin{lemma}[Finite-interval virial identities]\label{lem:virial}
Suppose $w$ solves~\eqref{eq:scaled} on a time interval $I$ and
$w,xw,|x|^2w\in C(I;H^k)$ for every integer $k\ge0$.
Then $M,E$ are constant and
\begin{align}
 V'&=-4\beta D+\frac{16}{3}W+\frac43S^2G(S,\tau),\label{eq:Vprime}\\
 D'&=-2\beta\|\nabla w\|_2^2-\frac43P-\frac23SG(S,\tau)
     =-4E-\frac23P-\frac23SG(S,\tau).\label{eq:Dprime}
\end{align}
\end{lemma}

\begin{proof}
Unitarity gives $\langle w,Q_S(w)\rangle=P\in\R$, which proves mass
conservation. In the real $L^2$ pairing, the derivative of $E$ is
$-\beta\Delta w+Q_S(w)$. Equation~\eqref{eq:scaled} reads
$w_\tau=i(-\beta\Delta w+Q_S(w))$, so $E'=0$.

For the variance calculation set $h=T_\sigma w$. The position commutator gives
\[
 T_\sigma|x|^2w=(|x|^2+4i\sigma x\cdot\nabla+6i\sigma-4\sigma^2\Delta)h.
\]
The nonlinear part of $V'$ is $-2$ times the imaginary part of the pairing
of this expression with $|h|^4h$. Integration by parts gives
\[
 \operatorname{Re}\int x\cdot\nabla\bar h\,|h|^4h=-G/2,
 \qquad \partial_\sigma G=6\operatorname{Im}\int\Delta\bar h\,|h|^4h.
\]
Its integrand is therefore $8\sigma G+(4/3)\sigma^2\partial_\sigma G$.
After integration in $\sigma$ this becomes
$(16/3)W+(4/3)S^2G(S,\tau)$. The linear contribution is $-4\beta D$,
proving~\eqref{eq:Vprime}.

For~\eqref{eq:Dprime}, use $w_\lambda(x)=\lambda^{3/2}w(\lambda x)$.
Then
\[
 P(w_\lambda)=\lambda^4\int_0^{\lambda^2S}G(\sigma,\tau)\,d\sigma,
 \qquad \left.\frac{d}{d\lambda}P(w_\lambda)\right|_{\lambda=1}
       =4P+2SG(S,\tau).
\]
The generator $\mathcal A=x\cdot\nabla+3/2$ is skew-adjoint, and
$D=\operatorname{Im}\langle w,\mathcal Aw\rangle$. Consequently
\[
 D'=2\operatorname{Im}\langle w_\tau,\mathcal Aw\rangle
    =-2\left.\frac{d}{d\lambda}E(w_\lambda)\right|_{\lambda=1}.
\]
Differentiating the kinetic term and the preceding formula for $P$ gives
\eqref{eq:Dprime}. The assumed weighted regularity justifies these pairings,
dilations and finite-interval integrations by parts.
\end{proof}

\section{Proof of norm inflation}\label{sec:proof}

Let $\phi(y)=e^{-|y|^2/4}$, and, for $g\ge2$, choose
\begin{equation}\label{eq:parameters}
 L=g^{-9/4},\qquad A=g^{11/4},\qquad u_{g,0}(x)=A\phi(x/L).
\end{equation}
Let $u_g$ be the smooth solution given by Lemma~\ref{lem:smooth}. Define
\[
 w_g(\tau,y)=A^{-1}u_g(\tau/(A^4L^2),Ly).
\]
For a scaled function, $T_\sigma[A f(\cdot/L)](x)=A(T_{\sigma/L^2}f)(x/L)$.
Changing the integration variable gives
\[
 Q_1[A f(\cdot/L)](x)=A^5L^2(Q_{L^{-2}}f)(x/L).
\]
Thus $w_g$ satisfies~\eqref{eq:scaled} exactly, with
\begin{equation}\label{eq:scales}
 \beta=(AL)^{-4}=g^{-2},\qquad S=L^{-2}=g^{9/2},
 \qquad \tau=g^{13/2}t,\qquad w_g(0)=\phi.
\end{equation}
The averaging interval remains finite. Fourier scaling also gives
\begin{equation}\label{eq:normscale}
 \|u_g(t)\|_{\dot H^{1/4}}
 =AL^{5/4}\|w_g(\tau)\|_{\dot H^{1/4}}
 =g^{-1/16}\|w_g(\tau)\|_{\dot H^{1/4}}.
\end{equation}

We use the variance and energy of $w_g$ on $[-g,0]$. Suppress its index, and set
$M_0=\|\phi\|_2^2>0$ and $V_0=\|y\phi\|_2^2$. The free Gaussian satisfies
\[
 T_\sigma\phi=(1+i\sigma)^{-3/2}
    \exp\left(-\frac{|y|^2}{4(1+i\sigma)}\right),
 \qquad
 \|T_\sigma\phi\|_6^6=(2\pi/3)^{3/2}(1+\sigma^2)^{-3}.
\]
Since $\int_0^\infty(1+\sigma^2)^{-3}d\sigma=3\pi/16$, the initial energy,
and hence the conserved energy, is bounded independently of $g$ by
\begin{equation}\label{eq:Estar}
 E\le E_*:=\tfrac12\|\nabla\phi\|_2^2+\frac{\pi}{32}(2\pi/3)^{3/2}.
\end{equation}

The Gaussian is real, so $D(0)=0$. Since $0\le P\le6E$,
\eqref{eq:Dprime} gives $D'\ge-8E-(2/3)SG(S,\tau)$. For $\tau\le r\le0$,
\[
 D(r)\le-8Er+\frac23S\int_r^0G(S,q)\,dq.
\]
Integrating~\eqref{eq:Vprime} backwards gives
\[
 V(\tau)=V_0+4\beta\int_\tau^0D(r)\,dr
 -\int_\tau^0\left(\frac{16}{3}W(r)+\frac43S^2G(S,r)\right)dr.
\]
The double integral of $G(S,q)$ equals $\int_\tau^0(q-\tau)G(S,q)dq$.
Combining these formulas yields
\begin{align}\label{eq:backward}
 V(\tau)\le {}&V_0+16\beta E\tau^2-\frac{16}{3}\int_\tau^0W(r)\,dr\nonumber\\
 &+\frac43S\int_\tau^0[2\beta(r-\tau)-S]G(S,r)\,dr.
\end{align}
For $-g\le\tau\le0$, we have $2\beta g\le S$ and $\beta\tau^2\le1$.
Both integral terms on the right are nonpositive. Therefore
\begin{equation}\label{eq:Vstar}
 V(\tau)\le V_*:=V_0+16E_*,\qquad -g\le\tau\le0.
\end{equation}
Energy and Cauchy--Schwarz now give
\[
 \|\nabla w(\tau)\|_2\le\sqrt{2E_*}\,g,
 \qquad |D(\tau)|\le C_Dg,\qquad C_D=\sqrt{2E_*V_*}.
\]
In the integral of~\eqref{eq:Vprime} from $-g$ to zero the endpoint
term is nonnegative. Using $V(-g)\ge0$, we obtain the uniform bound
\begin{equation}\label{eq:Wupper}
 \frac{16}{3}\int_{-g}^0W(r)\,dr
 \le V_0+4\beta\int_{-g}^0D(r)\,dr
 \le V_0+4C_D=:\frac{16}{3}C_W.
\end{equation}

Let $B_g=\max_{-g\le\tau\le0}\|w_g(\tau)\|_{\dot H^{1/4}}$.
This maximum is finite and attained by $H^1$ continuity. Lemma~\ref{lem:localization}
applies at every time with the same mass $M_0$, position bound $\sqrt{V_*}$,
and $s=1/4$. Its constant is independent of $g$, and it gives
\[
 cg(1+B_g)^{-8}\le\int_{-g}^0W(r)\,dr\le C_W.
\]
Thus, for all sufficiently large $g$,
\begin{equation}\label{eq:Bg}
 B_g\ge(cg/C_W)^{1/8}-1\ge c_1g^{1/8},
\end{equation}
where $c_1>0$ is fixed. A maximizing time $\tau_g$ is different from zero
once this lower bound exceeds $\|\phi\|_{\dot H^{1/4}}$. Set
$t_g=\tau_g g^{-13/2}<0$. Equations~\eqref{eq:normscale} and~\eqref{eq:Bg}
then give
\[
 \|u_{g,0}\|_{\dot H^{1/4}}=g^{-1/16}\|\phi\|_{\dot H^{1/4}},
 \qquad\|u_g(t_g)\|_{\dot H^{1/4}}\ge c_1g^{1/16},
 \qquad 0<|t_g|\le g^{-11/2}.
\]
Choosing $g$ large enough makes the first norm less than $\varepsilon$, the
second larger than $\varepsilon^{-1}$, and the nonzero time less than
$\varepsilon$ in absolute value. This proves Theorem~\ref{thm:inflation}.

The construction can be recorded without the rescaled quantities.
\begin{corollary}[Gaussian construction]\label{cor:construction}
There are constants $c>0$ and $g_0\ge2$ such that, for every real $g\ge g_0$,
the global smooth solution of~\eqref{eq:GT} with
\[
 u_{g,0}(x)=g^{11/4}\exp(-g^{9/2}|x|^2/4),\qquad x\in\R^3,
\]
has a time $-g^{-11/2}\le t_g<0$ satisfying
\[
 \|u_{g,0}\|_{\dot H^{1/4}}=C_\phi g^{-1/16},\qquad
 \|u_g(t_g)\|_{\dot H^{1/4}}\ge cg^{1/16}.
\]
Here $C_\phi=\|e^{-|\cdot|^2/4}\|_{\dot H^{1/4}}$, $u_g(0)=u_{g,0}$,
and $c,g_0$ depend only on the fixed Gaussian and the fixed cutoff in
Lemma~\ref{lem:localization}.
\end{corollary}

\paragraph{AI assistance.}
OpenAI Codex assisted the exploration of the argument, algebraic checks,
drafting, and manuscript review. The author is responsible for the mathematical
claims and the final text.

\begin{thebibliography}{9}
\bibitem{KM}
Jumpei Kawakami and Jason Murphy,
\emph{Small and large data scattering for the dispersion-managed NLS},
arXiv:2407.11151v1 (2024),
\url{https://arxiv.org/abs/2407.11151v1}.

\bibitem{KT}
Markus Keel and Terence Tao,
\emph{Endpoint Strichartz estimates},
American Journal of Mathematics \textbf{120} (1998), 955--980.

\bibitem{K}
Matthew Kowalski,
\emph{On ill-posedness for the Gabitov--Turitsyn equation},
arXiv:2510.11887v1 (2025),
\url{https://arxiv.org/abs/2510.11887v1}.

\bibitem{N}
L. Nirenberg,
\emph{On elliptic partial differential equations},
Annali della Scuola Normale Superiore di Pisa, Classe di Scienze,
Series 3, \textbf{13} (1959), no.~2, 115--162.
\end{thebibliography}
\end{document}
